% GENERATED FILE. Edit canonical lessons, then run npm run generate:books.
% canonical-source-hash: fnv1a64:3ed1e644ecaa1f89
% canonical-lessons: TE-C301-nighantuvu, TE-C301-bandi, TE-C301-ato, TE-C301-taksi, TE-C301-petrolu

\chapter{Dictionary, Cart, Auto-rickshaw, Taxi, Petrol}
\label{ch:te-dictionary-cart-auto-rickshaw-taxi-petrol}
\hlchaptermodality{301}

\begin{chapteropening}
\textbf{By the end of this chapter:} \emph{I can say a dictionary, a cart, an auto-rickshaw, a taxi and petrol.}

Complete the last lesson of chapter 301: I can say a dictionary, a cart, an auto-rickshaw, a taxi and petrol.
\end{chapteropening}

\section[\texorpdfstring{\te{నిఘంటువు} (nighaṇṭuvu)}{నిఘంటువు (nighaṇṭuvu)}]{\te{నిఘంటువు} (nighaṇṭuvu) — a dictionary}

\label{lesson:TE-C301-nighantuvu}



\begin{quote}
Before the new one: say the Telugu for a button, then the Telugu for a handkerchief.
\end{quote}

\subsection*{You'll want to know: \te{నిఘంటువు}}
\textbf{\te{నిఘంటువు}} — \emph{nighaṇṭuvu} — \textquotedblleft{}a dictionary\textquotedblright{}.

\begin{grammarlens}[title={What you've built}]
One more word for this chapter.
\end{grammarlens}

\subsection*{Your turn}
\begin{itemize}
\raggedright
  \item \emph{Say it:} \emph{nighaṇṭuvu}
  \item \emph{Say it:} \emph{nighaṇṭuvu}, once more
  \item \emph{Say it:} say \emph{nighaṇṭuvu}
  \item \emph{Recall:} say the Telugu for a button, then the Telugu for a handkerchief, then say \emph{nighaṇṭuvu} again
\end{itemize}

\begin{tcolorbox}[breakable,colback=teal!4,colframe=teal!35!black,title={Before you move on}]
What does \te{నిఘంటువు} mean? (\textquotedblleft{}A dictionary\textquotedblright{}.) Say it once more.
\end{tcolorbox}
\section[\texorpdfstring{\te{బండి} (baṇḍi)}{బండి (baṇḍi)}]{\te{బండి} (baṇḍi) — a cart; a vehicle}

\label{lesson:TE-C301-bandi}



\begin{quote}
Before the new one: say the Telugu for a handkerchief, then the Telugu for a dictionary.
\end{quote}

\subsection*{You'll want to know: \te{బండి}}
\textbf{\te{బండి}} — \emph{baṇḍi} — \textquotedblleft{}a cart; a vehicle\textquotedblright{}.

\begin{grammarlens}[title={What you've built}]
One more word for this chapter.
\end{grammarlens}

\subsection*{Your turn}
\begin{itemize}
\raggedright
  \item \emph{Say it:} \emph{baṇḍi}
  \item \emph{Say it:} \emph{baṇḍi}, once more
  \item \emph{Say it:} say \emph{baṇḍi}
  \item \emph{Recall:} say the Telugu for a handkerchief, then the Telugu for a dictionary, then say \emph{baṇḍi} again
\end{itemize}

\begin{tcolorbox}[breakable,colback=teal!4,colframe=teal!35!black,title={Before you move on}]
What does \te{బండి} mean? (\textquotedblleft{}A cart; a vehicle\textquotedblright{}.) Say it once more.
\end{tcolorbox}
\section[\texorpdfstring{\te{ఆటో} (āṭō)}{ఆటో (āṭō)}]{\te{ఆటో} (āṭō) — an auto-rickshaw}

\label{lesson:TE-C301-ato}



\begin{quote}
Before the new one: say the Telugu for a dictionary, then the Telugu for a cart.
\end{quote}

\subsection*{You'll want to know: \te{ఆటో}}
\textbf{\te{ఆటో}} — \emph{āṭō} — \textquotedblleft{}an auto-rickshaw\textquotedblright{}.

\begin{grammarlens}[title={What you've built}]
One more word for this chapter.
\end{grammarlens}

\subsection*{Your turn}
\begin{itemize}
\raggedright
  \item \emph{Say it:} \emph{āṭō}
  \item \emph{Say it:} \emph{āṭō}, once more
  \item \emph{Say it:} say \emph{āṭō}
  \item \emph{Recall:} say the Telugu for a dictionary, then the Telugu for a cart, then say \emph{āṭō} again
\end{itemize}

\begin{tcolorbox}[breakable,colback=teal!4,colframe=teal!35!black,title={Before you move on}]
What does \te{ఆటో} mean? (\textquotedblleft{}An auto-rickshaw\textquotedblright{}.) Say it once more.
\end{tcolorbox}
\section[\texorpdfstring{\te{టాక్సీ} (ṭāksī)}{టాక్సీ (ṭāksī)}]{\te{టాక్సీ} (ṭāksī) — a taxi}

\label{lesson:TE-C301-taksi}



\begin{quote}
Before the new one: say the Telugu for a cart, then the Telugu for an auto-rickshaw.
\end{quote}

\subsection*{You'll want to know: \te{టాక్సీ}}
\textbf{\te{టాక్సీ}} — \emph{ṭāksī} — \textquotedblleft{}a taxi\textquotedblright{}.

\begin{grammarlens}[title={What you've built}]
One more word for this chapter.
\end{grammarlens}

\subsection*{Your turn}
\begin{itemize}
\raggedright
  \item \emph{Say it:} \emph{ṭāksī}
  \item \emph{Say it:} \emph{ṭāksī}, once more
  \item \emph{Say it:} say \emph{ṭāksī}
  \item \emph{Recall:} say the Telugu for a cart, then the Telugu for an auto-rickshaw, then say \emph{ṭāksī} again
\end{itemize}

\begin{tcolorbox}[breakable,colback=teal!4,colframe=teal!35!black,title={Before you move on}]
What does \te{టాక్సీ} mean? (\textquotedblleft{}A taxi\textquotedblright{}.) Say it once more.
\end{tcolorbox}
\section[\texorpdfstring{\te{పెట్రోలు} (peṭrōlu)}{పెట్రోలు (peṭrōlu)}]{\te{పెట్రోలు} (peṭrōlu) — petrol}

\label{lesson:TE-C301-petrolu}



\begin{quote}
Before the new one: say the Telugu for an auto-rickshaw, then the Telugu for a taxi.
\end{quote}

\subsection*{You'll want to know: \te{పెట్రోలు}}
\textbf{\te{పెట్రోలు}} — \emph{peṭrōlu} — \textquotedblleft{}petrol\textquotedblright{}.

\begin{grammarlens}[title={What you've built}]
One more word for this chapter.
\end{grammarlens}

\subsection*{Your turn}
\begin{itemize}
\raggedright
  \item \emph{Say it:} \emph{peṭrōlu}
  \item \emph{Say it:} \emph{peṭrōlu}, once more
  \item \emph{Say it:} say \emph{peṭrōlu}
  \item \emph{Recall:} say the Telugu for an auto-rickshaw, then the Telugu for a taxi, then say \emph{peṭrōlu} again
\end{itemize}

\begin{tcolorbox}[breakable,colback=teal!4,colframe=teal!35!black,title={Before you move on}]
What does \te{పెట్రోలు} mean? (\textquotedblleft{}Petrol\textquotedblright{}.) Say it once more.
\end{tcolorbox}
